\documentclass[11pt,a4paper,twoside,openany]{report}

% ----------------------------
% Pacchetti
% ----------------------------
\usepackage[utf8]{inputenc}
\usepackage[T1]{fontenc}
\usepackage[italian]{babel}
\usepackage{lmodern}
\usepackage{setspace}
\usepackage{geometry}
\usepackage{fancyhdr}
\usepackage{titlesec}
\usepackage{tocloft}
\usepackage{csquotes}
\usepackage{graphicx}
\graphicspath{
  {./img/}
  {./img/ble/}
  {./img/ei/}
}
\usepackage{float}
\usepackage{caption}
\usepackage{subcaption}
\usepackage{xcolor}
\usepackage{listings}
\usepackage{amsmath}
\usepackage{amssymb}
\usepackage{siunitx}
\sisetup{output-decimal-marker={,}, group-separator={\,}, group-minimum-digits=4}
\usepackage{booktabs}
\usepackage{array}
\usepackage{enumitem}
\usepackage{mdframed}
\usepackage{tabularx}
\usepackage{multirow}
\usepackage{tikz}
\usetikzlibrary{arrows.meta, positioning, shapes.geometric, fit, calc}
\usepackage{pgfplots}
\pgfplotsset{compat=1.17}
\usepackage{hyperref}
\usepackage[all]{hypcap}

% ----------------------------
% Margini e spaziatura
% ----------------------------
\geometry{a4paper, top=2.5cm, bottom=2.5cm, left=2.7cm, right=2.5cm}
\setstretch{1.05}
\setlist{itemsep=2pt, topsep=4pt}
\setlength{\parindent}{1.5em}
\setlength{\parskip}{0pt}

% ----------------------------
% Stile capitoli e paragrafi
% ----------------------------
\titleclass{\chapter}{straight}
\titleformat{\chapter}[block]{\normalfont\LARGE\bfseries}{\thechapter.}{1em}{}
\titlespacing*{\chapter}{0pt}{26pt}{12pt}
\renewcommand{\thechapter}{\arabic{chapter}}
\titleformat{\section}{\normalfont\large\bfseries}{\thesection}{1em}{}
\titlespacing*{\section}{0pt}{18pt}{6pt}
\titleformat{\subsection}{\normalfont\normalsize\bfseries}{\thesubsection}{1em}{}
\titlespacing*{\subsection}{0pt}{12pt}{4pt}
\titleformat{\subsubsection}{\normalfont\normalsize\itshape}{\thesubsubsection}{1em}{}
\titlespacing*{\subsubsection}{0pt}{8pt}{2pt}

% ----------------------------
% Numerazione pagine
% ----------------------------
\pagestyle{fancy}
\fancyhf{}
\fancyhead{}
\fancyfoot[LE]{\thepage}
\fancyfoot[RO]{\thepage}
\fancyfoot[C]{}
\fancypagestyle{plain}{%
  \fancyhf{}
  \fancyhead{}
  \fancyfoot[LE]{\thepage}
  \fancyfoot[RO]{\thepage}
  \fancyfoot[C]{}
}
\renewcommand{\headrulewidth}{0pt}
\renewcommand{\footrulewidth}{0pt}

% ----------------------------
% Stile codice
% ----------------------------
\definecolor{codegreen}{rgb}{0,0.6,0}
\definecolor{codegray}{rgb}{0.5,0.5,0.5}
\definecolor{codepurple}{rgb}{0.58,0,0.82}
\definecolor{backcolour}{rgb}{0.95,0.95,0.92}

\lstdefinestyle{mystyle}{
    language=C++,
    backgroundcolor=\color{backcolour},
    commentstyle=\color{codegreen},
    keywordstyle=\color{magenta},
    numberstyle=\tiny\color{codegray},
    stringstyle=\color{codepurple},
    basicstyle=\ttfamily\footnotesize,
    breakatwhitespace=false,
    breaklines=true,
    captionpos=b,
    keepspaces=true,
    numbers=left,
    numbersep=5pt,
    showspaces=false,
    showstringspaces=false,
    showtabs=false,
    tabsize=2,
    literate={à}{{\`a}}1 {è}{{\`e}}1 {é}{{\'e}}1 {ì}{{\`i}}1 {ò}{{\`o}}1 {ù}{{\`u}}1
}
\lstset{style=mystyle}
\renewcommand{\lstlistingname}{Listato}
\renewcommand{\lstlistlistingname}{Elenco dei listati}

% ----------------------------
% Indice nero
% ----------------------------
\hypersetup{
  hidelinks,
  pdftitle={Smart Laundry Color Classifier},
  pdfauthor={Andrea Vaerini},
}

% ----------------------------
% Riquadro per note e criticità
% ----------------------------
\definecolor{ucborder}{RGB}{60,90,150}
\definecolor{ucbg}{RGB}{245,247,252}

\newenvironment{ucbox}{%
  \begin{mdframed}[
    backgroundcolor=ucbg,
    linecolor=ucborder,
    linewidth=1pt,
    innerleftmargin=10pt,
    innerrightmargin=10pt,
    innertopmargin=8pt,
    innerbottommargin=8pt,
    roundcorner=4pt
  ]
}{%
  \end{mdframed}
}

% ----------------------------
% Elenchi stilizzati
% ----------------------------
\newenvironment{toollist}{%
  \begin{description}[
    style=nextline,
    leftmargin=0pt,
    labelindent=0pt,
    font=\normalfont\bfseries,
    itemsep=6pt
  ]
}{%
  \end{description}
}

% ----------------------------
% Comandi di comodo
% ----------------------------
\newcommand{\rn}{r_{\mathrm{norm}}}
\newcommand{\gn}{g_{\mathrm{norm}}}
\newcommand{\bn}{b_{\mathrm{norm}}}
\newcommand{\uuid}[1]{\texttt{\small #1}}

% ============================================================
\begin{document}

% ----------------------------
% Frontespizio
% ----------------------------
\begin{titlepage}
\thispagestyle{empty}
    \centering
    \includegraphics[width=0.26\textwidth, height=3.5cm]{logo_uni.jpg}

    {\Large UNIVERSITÀ DEGLI STUDI DI BERGAMO \\[0.3cm]}
    {\large Scuola di Ingegneria \\[0.1cm]}
    {\large Corso di Laurea in Ingegneria Informatica \\[1.5cm]}
    {\large \textbf{LABORATORIO DI SENSORI} \\[1.5cm]}

    {\LARGE\bfseries SMART LAUNDRY COLOR CLASSIFIER\\[0.5cm]}
    {\large Sistema Embedded TinyML su Arduino Nano 33 BLE\\
    per la Classificazione Cromatica e Notifica Wireless \\[2cm]}

    \begin{minipage}[t]{0.45\textwidth}
        \raggedright
        \textbf{Candidato:} \\
        Andrea Vaerini \textbf{~1067398}
    \end{minipage}
    \hfill
    \begin{minipage}[t]{0.45\textwidth}
       \raggedleft
        \textbf{Docenti:} \\
        Prof. Gianluca Traversi \\
        Dott. Paolo Lazzaroni \\[0.5cm]
    \end{minipage}

    \vfill
    {Anno Accademico 2026/2027}
\end{titlepage}

\pagenumbering{arabic}

\thispagestyle{fancy}
\tableofcontents
\listoffigures
\listoftables
\newpage

% ============================================================
\chapter{Introduzione e motivazione}
\label{ch:intro}

\section{Automazione nello \textit{Smart Laundry}}

La separazione dei capi per colore prima del lavaggio è un'operazione
quotidiana tanto banale quanto soggetta a errore: un singolo indumento
rosso o blu inserito in un carico di bianchi è sufficiente a compromettere
l'intero lavaggio per effetto del rilascio di coloranti (\textit{color
bleeding}). Le lavatrici di fascia alta integrano già sensori di carico,
di torbidità dell'acqua e di schiuma, ma il riconoscimento del colore dei
capi resta affidato all'utente. Un sensore cromatico posto in prossimità
dell'oblò, capace di avvisare l'utente direttamente sullo smartphone nel
momento in cui inserisce un capo \enquote{incompatibile} con il programma
selezionato, rappresenta un'estensione naturale del paradigma
\textit{Smart Home}.

\section{Limiti dei sensori cromatici tradizionali}

I classificatori di colore più semplici confrontano le tre componenti
$R$, $G$, $B$ fornite da un sensore con soglie fisse determinate
empiricamente. Questo approccio presenta tre limiti principali:

\begin{itemize}
  \item \textbf{Dipendenza dall'illuminazione}: l'intensità assoluta
        dei canali varia di un ordine di grandezza passando da luce
        diretta a ombra, per cui una soglia tarata in una condizione non
        è valida nelle altre;
  \item \textbf{Confini di decisione non lineari}: colori come rosa,
        viola e bianco occupano regioni dello spazio cromatico contigue
        e di forma irregolare, difficili da separare con semplici
        disuguaglianze;
  \item \textbf{Ambiguità del nero}: un tessuto nero opaco e l'assenza
        di qualunque oggetto davanti al sensore producono entrambi
        letture prossime allo zero su tutti i canali e risultano quindi
        indistinguibili sulla sola base dell'informazione cromatica.
\end{itemize}

Il progetto affronta il primo punto con una normalizzazione cromatica
(Capitolo~\ref{ch:software}), il secondo sostituendo le soglie fisse con
una piccola rete neurale addestrata (\textit{TinyML}) e il terzo con un
\textit{Proximity Gate} basato sul sensore di prossimità a infrarossi
integrato (Capitolo~\ref{ch:analisi}).

\section{Requisiti del progetto}

Il sistema sviluppato, denominato \textit{Smart Laundry Color
Classifier}, deve soddisfare i seguenti requisiti:

\begin{toollist}
  \item[R1 -- Classificazione multi-classe]
        Riconoscere otto classi cromatiche rappresentative dei capi
        domestici: \texttt{BLACK}, \texttt{WHITE}, \texttt{RED},
        \texttt{BLUE}, \texttt{GREEN}, \texttt{PINK}, \texttt{PURPLE},
        \texttt{YELLOW}.
  \item[R2 -- Robustezza all'illuminazione]
        Fornire la stessa classificazione in luce diretta, in luce
        ambiente e in ombra.
  \item[R3 -- Rilevamento della presenza]
        Distinguere un capo nero dall'assenza di capi, evitando falsi
        positivi quando il sensore è scoperto.
  \item[R4 -- Inferenza a bordo]
        Eseguire l'intera elaborazione sul microcontrollore, senza
        dipendere da servizi cloud.
  \item[R5 -- Notifica wireless]
        Inviare il risultato (classe e confidenza) a uno smartphone
        tramite Bluetooth Low Energy.
\end{toollist}

\section{Sviluppo su singola scheda}

Il progetto è stato svolto individualmente. Per contenere i tempi di
sviluppo e massimizzare l'affidabilità meccanica del prototipo si è
scelto di utilizzare esclusivamente la sensoristica già presente sulla
scheda \textbf{Arduino Nano 33 BLE Sense}: il sensore di colore,
prossimità e gesti APDS-9960 e il modulo radio BLE sono integrati sul PCB
e collegati internamente al microcontrollore. Questa scelta elimina
completamente cablaggi su breadboard, resistenze di pull-up esterne e
problemi di contatti intermittenti, e consente di concentrare il lavoro
sugli aspetti di elaborazione del segnale, di apprendimento automatico e
di comunicazione. Il prototipo risultante è costituito da un unico
modulo di \SI{45}{mm} $\times$ \SI{18}{mm} alimentato via USB.

% ============================================================
\chapter{Architettura hardware e sensoristica integrata}
\label{ch:hardware}

\section{Arduino Nano 33 BLE Sense}

La Arduino Nano 33 BLE Sense è una scheda di sviluppo in formato
\textit{Nano} costruita attorno al modulo u-blox NINA-B306, che ospita il
SoC \textbf{Nordic Semiconductor nRF52840}~\cite{nrf52840}. Le
caratteristiche rilevanti per il progetto sono riassunte in
Tabella~\ref{tab:nano33}.

\begin{table}[htbp]
\centering
\caption{Specifiche principali della Arduino Nano 33 BLE Sense.}
\label{tab:nano33}
\begin{tabular}{@{}ll@{}}
\toprule
\textbf{Caratteristica} & \textbf{Valore} \\
\midrule
SoC                     & Nordic nRF52840 \\
CPU                     & ARM Cortex-M4F, \SI{64}{MHz}, FPU singola precisione \\
Memoria Flash           & \SI{1}{MB} (\SI{983040}{B} disponibili per lo sketch) \\
Memoria SRAM            & \SI{256}{kB} \\
Radio                   & Bluetooth 5 / BLE, \SI{2.4}{GHz} \\
Tensione di lavoro      & \SI{3.3}{V} (alimentazione USB \SI{5}{V}) \\
Bus sensori             & I\textsuperscript{2}C interno (\texttt{Wire1}) \\
Sensori integrati       & APDS-9960, IMU, microfono, temp./umidità, pressione \\
\bottomrule
\end{tabular}
\end{table}

La presenza dell'unità in virgola mobile (FPU) del Cortex-M4F è
particolarmente importante: permette di eseguire sia la normalizzazione
cromatica sia l'inferenza della rete neurale direttamente in
\texttt{float32}, senza ricorrere alla quantizzazione a 8 bit e quindi
senza perdita di precisione numerica. Il firmware è basato sul core
\texttt{Arduino Mbed OS Nano Boards} (versione 1.1.6), che fornisce lo
stack BLE Cordio utilizzato dalla libreria \texttt{ArduinoBLE}.

\section{Il sensore Avago/Broadcom APDS-9960}

L'APDS-9960~\cite{apds9960} è un sensore digitale che integra in un
unico package ottico quattro funzioni: misura di luce ambiente e colore
(ALS/RGBC), rilevamento di prossimità, riconoscimento di gesti e un
emettitore IR. Comunica tramite I\textsuperscript{2}C all'indirizzo
fisso \texttt{0x39} (registro ID atteso: \texttt{0xAB}, verificato dal
firmware in fase di diagnostica).

\begin{figure}[htbp]
\centering
\begin{tikzpicture}[
  node distance=0.6cm and 1.0cm,
  blk/.style={draw, rounded corners, minimum width=2.6cm, minimum height=0.9cm, align=center, font=\small},
  sens/.style={blk, fill=ucbg},
  arr/.style={-{Stealth[length=2.5mm]}, thick}
]
  \node[sens] (pd) {Fotodiodi\\RGBC + filtri};
  \node[sens, below=of pd] (ir) {Fotodiodo IR\\(prossimità)};
  \node[sens, below=of ir] (led) {LED IR\\\SI{950}{nm}};
  \node[blk, right=of pd] (adc) {ADC 16 bit\\(ATIME, AGAIN)};
  \node[blk, right=of ir] (padc) {ADC 8 bit\\(PDATA)};
  \node[blk, right=1.6cm of $(adc)!0.5!(padc)$] (i2c) {Interfaccia\\I\textsuperscript{2}C \texttt{0x39}};
  \node[blk, right=of i2c, fill=ucbg] (mcu) {nRF52840};
  \node[draw, dashed, inner sep=6pt, fit=(pd)(led)(adc)(padc)(i2c), label={[font=\footnotesize]above:APDS-9960}] {};
  \draw[arr] (pd) -- (adc);
  \draw[arr] (ir) -- (padc);
  \draw[arr] (adc) -- (i2c);
  \draw[arr] (padc) -- (i2c);
  \draw[arr] (i2c) -- (mcu);
  \draw[arr, dashed] (led.east) -| node[pos=0.75, right, font=\footnotesize, align=left]{eco IR\\dal tessuto} (padc.south);
\end{tikzpicture}
\caption{Schema a blocchi funzionale dell'APDS-9960 e collegamento al SoC.}
\label{fig:apds_block}
\end{figure}

\subsection{Canali RGBC}

Il blocco colore è costituito da una matrice di fotodiodi ricoperti da
filtri interferenziali rosso, verde e blu, più un canale
\textit{Clear} senza filtro cromatico. Ciascun fotodiodo genera una
fotocorrente proporzionale all'irradianza incidente pesata per la
propria responsività spettrale $S_k(\lambda)$:
\begin{equation}
  I_k = \int_{\lambda} S_k(\lambda)\, E(\lambda)\, \rho(\lambda)\, d\lambda,
  \qquad k \in \{R, G, B, C\}
  \label{eq:fotocorrente}
\end{equation}
dove $E(\lambda)$ è la distribuzione spettrale della sorgente luminosa e
$\rho(\lambda)$ la riflettanza del tessuto. Le fotocorrenti sono
integrate da quattro ADC a doppia rampa da 16 bit. Due registri
determinano la sensibilità:

\begin{itemize}
  \item \textbf{ATIME} (tempo di integrazione): ogni ciclo dura
        \SI{2.78}{ms} e il valore massimo del conteggio vale
        $1025 \cdot (256 - \mathrm{ATIME})$;
  \item \textbf{AGAIN} (guadagno analogico): selezionabile tra
        $1\times$, $4\times$, $16\times$ e $64\times$.
\end{itemize}

La libreria \texttt{Arduino\_APDS9960} imposta in fase di
inizializzazione $\mathrm{ATIME} = 253$ (3 cicli, circa
\SI{8.3}{ms}, fondo scala 3075 conteggi) e $\mathrm{AGAIN} = 16\times$.
Come discusso nel Capitolo~\ref{ch:analisi}, questa configurazione,
pensata per letture rapide in luce ambiente, determina conteggi molto
bassi quando il tessuto è appoggiato sul sensore.

È importante sottolineare che l'APDS-9960 \textbf{non dispone di un
illuminatore nel visibile}: le letture di colore dipendono interamente
dalla luce ambiente riflessa dal tessuto, il che rende indispensabile la
normalizzazione descritta nel Capitolo~\ref{ch:software}.

\subsection{Rilevamento di prossimità}

Il blocco di prossimità sfrutta un LED IR integrato (\SI{950}{nm}) che
emette una serie di impulsi (configurati dalla libreria a 16 impulsi da
\SI{16}{\micro s}, registro \texttt{PPULSE = 0x8F}). La radiazione
riflessa da un oggetto vicino viene raccolta da un fotodiodo IR
dedicato; il segnale, depurato dalla componente ambiente, è convertito
in un valore a 8 bit \texttt{PDATA} che \textbf{cresce all'avvicinarsi
dell'oggetto}. La portata utile per superfici tessili è di pochi
centimetri, ed è proprio questa caratteristica a rendere il sensore
adatto a discriminare la presenza fisica di un capo indipendentemente
dal suo colore: anche un tessuto nero riflette una frazione non nulla
della radiazione nel vicino infrarosso.

% ============================================================
\chapter{Pipeline software e TinyML}
\label{ch:software}

\section{Normalizzazione cromatica}
\label{sec:normalizzazione}

Dall'Equazione~\eqref{eq:fotocorrente} si osserva che, se la sorgente
luminosa varia soltanto in intensità, ovvero $E(\lambda) \rightarrow
k\,E(\lambda)$, tutti i canali vengono moltiplicati per lo stesso fattore
$k$. Il rapporto tra un canale e la somma dei tre canali cromatici è
quindi indipendente da $k$. Su questo principio si basano le
\textbf{coordinate cromatiche normalizzate}, calcolate a partire dai
conteggi grezzi $R$, $G$, $B$ forniti dall'APDS-9960:
\begin{equation}
  \rn = \frac{R}{R+G+B}, \qquad
  \gn = \frac{G}{R+G+B}, \qquad
  \bn = \frac{B}{R+G+B}.
  \label{eq:norm}
\end{equation}

Per costruzione vale $\rn + \gn + \bn = 1$: le tre feature giacciono su
un piano (simplesso cromatico) e i gradi di libertà effettivi sono due.
L'informazione sulla luminosità complessiva, che varia con
l'illuminazione, viene scartata, mentre si conserva la
\textit{cromaticità}, legata alla riflettanza spettrale $\rho(\lambda)$
del tessuto. Il canale \textit{Clear} non entra nelle feature ma viene
registrato nelle misure come indicatore del livello di segnale.

Il caso degenere $R+G+B = 0$, che si presenta con tessuti neri o in
completa oscurità, è gestito ponendo il denominatore pari a 1: il punto
risultante $(0,0,0)$ si trova fuori dal simplesso ed è quindi
facilmente riconoscibile dal classificatore, ma è \emph{identico} per un
capo nero e per l'assenza di capi. Questo aspetto è ripreso nella
Sezione~\ref{sec:nero_vuoto}.

\section{Architettura del firmware}

Il firmware (\texttt{firmware/SmartLaundry/SmartLaundry.ino}) è scritto
in C++ per l'ambiente Arduino e si appoggia a tre librerie principali,
riassunte in Tabella~\ref{tab:librerie}.

\begin{table}[htbp]
\centering
\caption{Librerie software utilizzate nel firmware.}
\label{tab:librerie}
\begin{tabularx}{\textwidth}{@{}>{\raggedright\arraybackslash}p{4.3cm}X@{}}
\toprule
\textbf{Libreria} & \textbf{Ruolo} \\
\midrule
\texttt{Arduino\_APDS9960} & Inizializzazione del sensore, lettura dei
  canali RGBC (\texttt{readColor}) e della prossimità
  (\texttt{readProximity}) via I\textsuperscript{2}C. \\
\texttt{Smart-Laundry-Color-}\newline\texttt{Classifier\_inferencing} & Libreria C++
  generata da Edge Impulse: blocco DSP, rete neurale compilata con EON
  Compiler e funzione \texttt{run\_classifier()}. \\
\texttt{ArduinoBLE} & Gestione del ruolo \textit{peripheral} BLE:
  advertising, servizio GATT custom e notifiche. \\
\bottomrule
\end{tabularx}
\end{table}

\subsection{Modello TinyML con Edge Impulse}

L'addestramento è stato svolto sulla piattaforma Edge
Impulse~\cite{edgeimpulse}. Per ogni classe sono stati acquisiti
campioni delle tre coordinate normalizzate tramite porta seriale, nelle
diverse condizioni di illuminazione descritte nel
Capitolo~\ref{ch:analisi}. L'\textit{impulse} è composto da:

\begin{toollist}
  \item[Ingresso]
        Finestra di un singolo campione a tre assi
        ($\rn$, $\gn$, $\bn$), frequenza nominale \SI{10}{Hz}.
  \item[Blocco di elaborazione \textit{Flatten}]
        Calcola la media di ciascun asse sulla finestra (unica
        statistica abilitata, scala 1,0), producendo un vettore di 3
        feature.
  \item[Blocco di apprendimento \textit{Classification}]
        Rete neurale densa con un livello nascosto da 16 neuroni ReLU e
        uno strato di uscita da 8 neuroni con attivazione
        \textit{softmax}, per un totale di
        $(3 \cdot 16 + 16) + (16 \cdot 8 + 8) = 200$ parametri.
\end{toollist}

Il modello è stato esportato come libreria Arduino utilizzando l'\textbf{EON
Compiler}, che, invece di includere un interprete TensorFlow Lite
generico, genera direttamente codice C++ contenente il grafo, i pesi
come array costanti in Flash e un'allocazione statica dei tensori.
È stata scelta la variante \textbf{non ottimizzata (float32)}: data la
dimensione ridottissima della rete (\SI{800}{B} di pesi), la
quantizzazione \texttt{int8} non avrebbe portato benefici apprezzabili
in memoria, mentre avrebbe introdotto errori di quantizzazione su
feature già fortemente discretizzate (Sezione~\ref{sec:illuminazione}).
L'arena dei tensori generata dal compilatore occupa soltanto
\SI{336}{B} di RAM.

\begin{lstlisting}[caption={Normalizzazione cromatica e inferenza (estratto da \texttt{loop()}).}, label={lst:inferenza}]
int r_raw, g_raw, b_raw, c_raw;
APDS.readColor(r_raw, g_raw, b_raw, c_raw);

// Normalizzazione delle coordinate (r_norm, g_norm, b_norm)
float sum = (float)(r_raw + g_raw + b_raw);
if (sum == 0.0f) sum = 1.0f;

float r_norm = (float)r_raw / sum;
float g_norm = (float)g_raw / sum;
float b_norm = (float)b_raw / sum;

for (size_t ix = 0; ix < EI_CLASSIFIER_DSP_INPUT_FRAME_SIZE; ix += 3) {
    features[ix + 0] = r_norm;
    features[ix + 1] = g_norm;
    features[ix + 2] = b_norm;
}

signal_t features_signal;
features_signal.total_length = EI_CLASSIFIER_DSP_INPUT_FRAME_SIZE;
features_signal.get_data = &raw_feature_get_data;

ei_impulse_result_t result = { 0 };
EI_IMPULSE_ERROR res = run_classifier(&features_signal, &result, false);
\end{lstlisting}

Al termine dell'inferenza il firmware individua la classe con
probabilità massima (\textit{arg\,max} sull'uscita softmax) e costruisce
il payload testuale nel formato \texttt{"CLASSE (NN\%)"}, ad esempio
\texttt{"BLUE (94\%)"}.

\subsection{Stack Bluetooth Low Energy}

La scheda opera come \textit{peripheral} GATT con nome
\texttt{SmartLaundry-Nano33}. È stato definito un servizio custom
contenente un'unica caratteristica di tipo stringa (massimo 64 byte) con
proprietà \texttt{READ} e \texttt{NOTIFY}, così che lo smartphone riceva
automaticamente ogni nuovo risultato senza effettuare
\textit{polling}.

\begin{table}[htbp]
\centering
\caption{Profilo GATT esposto dal dispositivo.}
\label{tab:gatt}
\begin{tabular}{@{}lll@{}}
\toprule
\textbf{Elemento} & \textbf{UUID} & \textbf{Proprietà} \\
\midrule
Servizio        & \uuid{19B10000-E8F2-537E-4F6C-D104768A1214} & Primary \\
Caratteristica  & \uuid{19B10001-E8F2-537E-4F6C-D104768A1214} & Read, Notify \\
\bottomrule
\end{tabular}
\end{table}

\begin{lstlisting}[caption={Definizione del servizio GATT e avvio dell'advertising.}, label={lst:ble}]
BLEService laundryService("19B10000-E8F2-537E-4F6C-D104768A1214");
BLEStringCharacteristic colorResultChar(
    "19B10001-E8F2-537E-4F6C-D104768A1214", BLERead | BLENotify, 64);

BLE.setLocalName("SmartLaundry-Nano33");
BLE.setAdvertisedService(laundryService);
laundryService.addCharacteristic(colorResultChar);
BLE.addService(laundryService);
BLE.advertise();
\end{lstlisting}

I payload possibili sono il nome di una delle otto classi seguito dalla
confidenza percentuale, oppure il messaggio di stato
\texttt{NESSUN\_CAPO} generato dal Proximity Gate.

\section{Macchina a stati a due stadi}
\label{sec:fsm}

Il comportamento del firmware è descritto dalla macchina a stati di
Figura~\ref{fig:fsm}. Dopo l'inizializzazione, in cui il sensore viene
avviato con un massimo di cinque tentativi e, in caso di fallimento,
viene eseguita una scansione diagnostica del bus I\textsuperscript{2}C,
il ciclo principale (periodo di circa \SI{500}{ms}) è articolato in due
stadi:

\begin{enumerate}
  \item \textbf{Proximity Gate}: viene letto il valore di prossimità IR.
        Se nessun oggetto si trova entro circa 2--3\,cm dal sensore, il
        sistema resta in \textit{standby}, invia la notifica
        \texttt{NESSUN\_CAPO} e non esegue alcuna classificazione;
  \item \textbf{Classificazione TinyML}: solo se il capo è presente si
        acquisiscono i canali RGBC, si calcolano le coordinate
        normalizzate, si esegue l'inferenza e si notifica il risultato.
\end{enumerate}

\begin{figure}[htbp]
\centering
\begin{tikzpicture}[
  node distance=0.9cm and 1.3cm,
  st/.style={draw, rounded corners, fill=ucbg, minimum width=2.7cm, minimum height=1cm, align=center, font=\small},
  dec/.style={draw, diamond, aspect=2, fill=white, align=center, font=\small, inner sep=1pt},
  arr/.style={-{Stealth[length=2.5mm]}, thick}
]
  \node[st] (init) {INIT\\APDS + BLE};
  \node[dec, below=of init] (gate) {Capo\\presente?};
  \node[st, right=1.6cm of gate] (idle) {STANDBY\\notifica\\\texttt{NESSUN\_CAPO}};
  \node[st, below=of gate] (acq) {ACQUISIZIONE\\RGBC + norm.};
  \node[st, below=of acq] (inf) {INFERENZA\\EON (float32)};
  \node[st, right=1.6cm of inf] (notify) {NOTIFICA BLE\\\texttt{CLASSE (NN\%)}};
  \node[draw, dashed, rounded corners, left=1.2cm of acq, align=center, font=\footnotesize] (wait) {attesa\\\SI{500}{ms}};
  \draw[arr] (init) -- (gate);
  \draw[arr] (gate) -- node[above, font=\footnotesize]{no} (idle);
  \draw[arr] (gate) -- node[right, font=\footnotesize]{sì} (acq);
  \draw[arr] (acq) -- (inf);
  \draw[arr] (inf) -- (notify);
  \draw[arr] (idle.north) |- ++(0,0.6) -| (wait.north);
  \draw[arr] (notify.south) |- ++(0,-0.6) -| (wait.south);
  \draw[arr] (wait.east) -- ++(0.4,0) |- (gate.west);
\end{tikzpicture}
\caption{Macchina a stati del firmware: Proximity Gate seguito dalla classificazione.}
\label{fig:fsm}
\end{figure}

La scelta di porre il gate \emph{prima} dell'acquisizione cromatica ha
un duplice vantaggio: elimina alla radice la classificazione di
\enquote{capi fantasma} quando il sensore è scoperto e, in un'ottica di
prodotto alimentato a batteria, consente di evitare l'esecuzione della
rete neurale e la trasmissione di risultati privi di significato.

% ============================================================
\chapter{Analisi sperimentale e discussione delle criticità}
\label{ch:analisi}

\section{Metodologia di misura}

Per caratterizzare il comportamento del sensore è stato sviluppato uno
sketch di acquisizione che stampa su porta seriale, in formato CSV, i
conteggi grezzi $R$, $G$, $B$, $C$ e le coordinate normalizzate
dell'Equazione~\eqref{eq:norm}. Per ciascuno degli otto colori è stato
utilizzato un campione di tessuto appoggiato sulla finestra ottica
dell'APDS-9960, acquisendo da 9 a 14 letture consecutive in tre
condizioni di illuminazione:

\begin{itemize}
  \item \textbf{Luce diretta}: tessuto illuminato da una torcia/lampada
        puntata sul campione;
  \item \textbf{Luce ambiente}: normale illuminazione interna del
        laboratorio;
  \item \textbf{Ombra}: sensore e campione schermati dalla luce
        diretta.
\end{itemize}

I file grezzi (\texttt{report/data/<COLORE>\_<condizione>.txt}) sono
inclusi nel repository. La Tabella~\ref{tab:misure} riassume, per ogni
combinazione, il numero di letture $n$, la media del canale
\textit{Clear} $\bar{C}$, le medie delle tre coordinate normalizzate e la
deviazione standard campionaria $\sigma_r$ di $\rn$.

\begin{table}[htbp]
\centering
\caption{Statistiche delle letture per colore e condizione di illuminazione.}
\label{tab:misure}
\small
\begin{tabular}{@{}llrrcccc@{}}
\toprule
\textbf{Colore} & \textbf{Condizione} & $n$ & $\bar{C}$ & $\bar{r}_{\mathrm{norm}}$ & $\bar{g}_{\mathrm{norm}}$ & $\bar{b}_{\mathrm{norm}}$ & $\sigma_r$ \\
\midrule
\multirow{3}{*}{Nero}   & Luce diretta  & 10 & 0,7  & 0,000 & 0,000 & 0,000 & 0,000 \\
                        & Luce ambiente & 10 & 0,5  & 0,000 & 0,000 & 0,000 & 0,000 \\
                        & Ombra         & 10 & 0,0  & 0,000 & 0,000 & 0,000 & 0,000 \\
\midrule
\multirow{3}{*}{Bianco} & Luce diretta  & 10 & 25,3 & 0,530 & 0,258 & 0,212 & 0,009 \\
                        & Luce ambiente & 11 & 3,3  & 0,539 & 0,261 & 0,200 & 0,070 \\
                        & Ombra         & 12 & 4,2  & 0,531 & 0,240 & 0,229 & 0,046 \\
\midrule
\multirow{3}{*}{Rosso}  & Luce diretta  & 9  & 3,8  & 0,596 & 0,280 & 0,123 & 0,067 \\
                        & Luce ambiente & 11 & 17,4 & 0,546 & 0,265 & 0,189 & 0,008 \\
                        & Ombra         & 12 & 1,8  & 0,611 & 0,389 & 0,000 & 0,148 \\
\midrule
\multirow{3}{*}{Blu}    & Luce diretta  & 10 & 2,7  & 0,366 & 0,316 & 0,316 & 0,070 \\
                        & Luce ambiente & 11 & 3,9  & 0,609 & 0,286 & 0,105 & 0,075 \\
                        & Ombra         & 12 & 6,9  & 0,570 & 0,234 & 0,196 & 0,044 \\
\midrule
\multirow{3}{*}{Verde}  & Luce diretta  & 12 & 27,8 & 0,431 & 0,309 & 0,260 & 0,045 \\
                        & Luce ambiente & 11 & 12,5 & 0,469 & 0,291 & 0,240 & 0,066 \\
                        & Ombra         & 11 & 13,0 & 0,478 & 0,267 & 0,255 & 0,029 \\
\midrule
\multirow{3}{*}{Rosa}   & Luce diretta  & 14 & 8,0  & 0,725 & 0,123 & 0,152 & 0,087 \\
                        & Luce ambiente & 11 & 2,3  & 0,964 & 0,018 & 0,018 & 0,121 \\
                        & Ombra         & 11 & 2,6  & 0,927 & 0,036 & 0,036 & 0,162 \\
\midrule
\multirow{3}{*}{Viola}  & Luce diretta  & 13 & 8,0  & 0,524 & 0,296 & 0,181 & 0,096 \\
                        & Luce ambiente & 11 & 3,9  & 0,463 & 0,278 & 0,260 & 0,057 \\
                        & Ombra$^{*}$   & 11 & 3,9  & 0,463 & 0,278 & 0,260 & 0,057 \\
\midrule
\multirow{3}{*}{Giallo} & Luce diretta  & 11 & 5,2  & 0,592 & 0,281 & 0,127 & 0,062 \\
                        & Luce ambiente & 10 & 4,0  & 0,600 & 0,233 & 0,167 & 0,061 \\
                        & Ombra         & 10 & 4,4  & 0,583 & 0,225 & 0,192 & 0,063 \\
\bottomrule
\multicolumn{8}{@{}l}{\footnotesize $^{*}$ Serie identica a quella in luce ambiente: verosimile duplicazione in fase di salvataggio.}
\end{tabular}
\end{table}

\section{Stabilità delle letture al variare dell'illuminazione}
\label{sec:illuminazione}

\subsection{Efficacia della normalizzazione}

Il campione bianco costituisce il caso più chiaro: passando dalla luce
diretta alla luce ambiente il canale \textit{Clear} si riduce di quasi
un fattore 8 (da 25,3 a 3,3 conteggi), mentre $\bar{r}_{\mathrm{norm}}$
varia di meno di 0,01 (0,530 contro 0,539) e le altre due coordinate di
circa 0,02. Un comportamento analogo si osserva per il verde (variazione
di $\bar{C}$ superiore a un fattore 2, scostamento massimo di
$\bar{r}_{\mathrm{norm}}$ pari a 0,047) e per il giallo (scostamento
massimo 0,017). La normalizzazione rende quindi le feature largamente
indipendenti dall'intensità della luce, come previsto
dalla Sezione~\ref{sec:normalizzazione}.

\begin{figure}[htbp]
\centering
\begin{tikzpicture}
\begin{axis}[
  width=\textwidth, height=7cm,
  xlabel={$\bar{r}_{\mathrm{norm}}$}, ylabel={$\bar{g}_{\mathrm{norm}}$},
  xmin=0.33, xmax=1.0, ymin=0, ymax=0.42,
  grid=major, grid style={gray!25},
  legend style={font=\scriptsize, at={(0.98,0.98)}, anchor=north east, cells={anchor=west}},
  legend columns=2,
  every axis plot/.append style={only marks, mark size=2.6pt, thick},
  x tick label style={/pgf/number format/use comma}, y tick label style={/pgf/number format/use comma},
]
\addplot[mark=o, color=black!55]          coordinates {(0.530,0.258) (0.539,0.261) (0.531,0.240)};
\addlegendentry{Bianco}
\addplot[mark=square*, color=red!80!black] coordinates {(0.596,0.280) (0.546,0.265) (0.611,0.389)};
\addlegendentry{Rosso}
\addplot[mark=triangle*, color=blue!80]   coordinates {(0.366,0.316) (0.609,0.286) (0.570,0.234)};
\addlegendentry{Blu}
\addplot[mark=diamond*, color=green!55!black] coordinates {(0.431,0.309) (0.469,0.291) (0.478,0.267)};
\addlegendentry{Verde}
\addplot[mark=pentagon*, color=magenta!70] coordinates {(0.725,0.123) (0.964,0.018) (0.927,0.036)};
\addlegendentry{Rosa}
\addplot[mark=*, color=violet]            coordinates {(0.524,0.296) (0.463,0.278)};
\addlegendentry{Viola}
\addplot[mark=star, color=orange!90!black, mark size=3.2pt] coordinates {(0.592,0.281) (0.600,0.233) (0.583,0.225)};
\addlegendentry{Giallo}
\end{axis}
\end{tikzpicture}
\caption{Medie delle coordinate cromatiche nel piano $(\rn, \gn)$ per le tre
condizioni di illuminazione. Il nero, sempre in $(0,0)$, è fuori scala.}
\label{fig:piano_cromatico}
\end{figure}

\subsection{Limiti osservati: segnale basso e quantizzazione}

L'analisi della Tabella~\ref{tab:misure} e della
Figura~\ref{fig:piano_cromatico} evidenzia però anche alcune criticità
fisiche che è importante discutere.

\begin{toollist}
  \item[Conteggi molto bassi]
        Con il tessuto appoggiato sulla finestra ottica, il campione
        stesso scherma gran parte della luce ambiente e, non essendoci
        un illuminatore nel visibile, il segnale che raggiunge i
        fotodiodi è minimo. Con la configurazione della libreria
        (\SI{8.3}{ms}, guadagno $16\times$, fondo scala 3075 conteggi),
        $\bar{C}$ non supera 28 conteggi, cioè meno dell'1\,\% della
        dinamica disponibile.
  \item[Quantizzazione delle feature]
        Con somme $R+G+B$ dell'ordine di 3--5 conteggi, le coordinate
        normalizzate possono assumere soltanto pochi valori discreti
        (0,250; 0,333; 0,500; 0,600; 0,667\,\dots), come si vede
        chiaramente nei file grezzi. La risoluzione di $\rn$ è dunque
        dell'ordine di $1/(R+G+B) \approx 0{,}2$--$0{,}3$, e una singola
        unità di rumore sul canale blu sposta il punto di molti
        centesimi. Le serie con $\bar{C}$ elevato (bianco in luce
        diretta, rosso in luce ambiente) mostrano infatti $\sigma_r <
        0{,}01$, mentre quelle a basso segnale arrivano a
        $\sigma_r \approx 0{,}15$.
  \item[Sbilanciamento verso il rosso]
        Anche il bianco presenta $\rn \approx 0{,}53$ anziché il valore
        ideale $1/3$. Ciò è dovuto alla combinazione della
        responsività spettrale dei fotodiodi, maggiore alle lunghezze
        d'onda lunghe, e dello spettro \enquote{caldo} delle sorgenti
        utilizzate. Lo sbilanciamento è sistematico e quindi viene
        appreso dal classificatore, ma riduce la distanza tra le classi.
  \item[Classi sovrapposte]
        Bianco, viola, verde, giallo e rosso occupano una regione
        compatta del piano cromatico ($\rn \in [0{,}43;\,0{,}61]$,
        $\gn \in [0{,}22;\,0{,}31]$). Il blu in particolare è ben
        distinto solo in luce diretta ($\rn = 0{,}366$), mentre in luce
        ambiente e in ombra si avvicina al giallo. Al contrario rosa e
        nero sono separati nettamente da tutte le altre classi.
\end{toollist}

Queste osservazioni giustificano la scelta di un classificatore non
lineare addestrato su dati raccolti in tutte le condizioni di
illuminazione, anziché di semplici soglie, e anticipano i risultati
della Sezione~\ref{sec:collaudo}: confidenze elevate per le classi
isolate e confidenze più basse per quelle che condividono la regione
centrale del piano cromatico.

\section{Ambiguità \enquote{nero contro vuoto} e Proximity Gate}
\label{sec:nero_vuoto}

Le misure sul campione nero mostrano $R = G = B = 0$ in tutte le
letture e in tutte le condizioni di illuminazione: il tessuto nero
opaco assorbe pressoché tutta la radiazione visibile e le coordinate
normalizzate collassano nel punto $(0,0,0)$. La stessa lettura si
ottiene però quando davanti al sensore non c'è alcun oggetto e la scena
è poco illuminata (per esempio all'interno del cestello). Dal punto di
vista del classificatore cromatico, quindi, \enquote{nero} e
\enquote{nessun capo} sono la stessa classe: senza un'informazione
aggiuntiva il sistema notificherebbe continuamente \texttt{BLACK} a
sensore scoperto.

\begin{ucbox}
\textbf{Soluzione adottata.} L'informazione mancante è fornita dal
canale di prossimità IR dello stesso APDS-9960. Il LED a \SI{950}{nm}
illumina attivamente la scena e la quantità di radiazione riflessa
dipende principalmente dalla \emph{distanza} della superficie, non dal
suo colore nel visibile: un tessuto nero appoggiato al sensore produce
un valore \texttt{PDATA} elevato, mentre l'aria non restituisce alcun
eco. Il firmware considera presente un capo soltanto quando
\texttt{PDATA} supera una soglia corrispondente a circa 2--3\,cm; in
caso contrario resta in standby e notifica \texttt{NESSUN\_CAPO}.
\end{ucbox}

La soglia è stata fissata sperimentalmente osservando il valore di
prossimità stampato ogni \SI{2}{s} sulla porta seriale nel messaggio
di stato del firmware, mentre si avvicinavano progressivamente i
campioni. Va notato che la
libreria \texttt{Arduino\_APDS9960} restituisce da
\texttt{readProximity()} il valore complementare $255 -
\texttt{PDATA}$, crescente con la distanza; il confronto con la soglia
deve quindi essere impostato tenendo conto di questa inversione di
scala. Il gate a due stadi (Figura~\ref{fig:fsm}) separa così in modo
netto i due casi: i falsi positivi \texttt{BLACK} a sensore scoperto
vengono eliminati, come confermato dal collaudo BLE
(Sezione~\ref{sec:collaudo}).

% ==== FINE CAPITOLI ====

% ============================================================
\begin{thebibliography}{9}
\addcontentsline{toc}{chapter}{Bibliografia}

\bibitem{nrf52840}
Nordic Semiconductor,
\textit{nRF52840 Product Specification v1.7}, 2021.

\bibitem{apds9960}
Broadcom (Avago Technologies),
\textit{APDS-9960 Digital Proximity, Ambient Light, RGB and Gesture
Sensor -- Data Sheet}, AV02-4191EN, 2015.

\bibitem{edgeimpulse}
S.~Hymel \textit{et al.},
\textit{Edge Impulse: An MLOps Platform for Tiny Machine Learning},
Proceedings of Machine Learning and Systems (MLSys), 2023.

\bibitem{arduinoapds}
Arduino,
\textit{Arduino\_APDS9960 Library},
\url{https://github.com/arduino-libraries/Arduino_APDS9960}.

\end{thebibliography}

\end{document}
